%------ SETUP OF THE DOCUMENT ------%
%This part of the document sets up the document and makes it easier to read the main.tex file
%Changes in the setup file are recommended if you wish to customize things like colors in links and such.
%------ ********************* ------%
% \usepackage[margin = 25mm]{geometry} %
\usepackage[a4paper,top=2.5cm,bottom=2.5cm,outer=2.25cm,inner=3cm]{geometry} %,parskip=1.5
\usepackage[T1]{fontenc}
\usepackage{helvet}
\usepackage{graphicx} % Required for inserting images
\usepackage[toc,page]{appendix}
\usepackage[english=usenglishmax]{hyphsubst} %sets hyphenation to American English. Google is your friend on hyphenation. Babel could also be used
\usepackage[hyphens]{url}
\usepackage{amsmath, enumitem, amssymb, array, setspace, float, subcaption, caption, booktabs, pdflscape, dcolumn, titlesec, tocloft, comment, xcolor, longtable, blindtext, rotating, lipsum, wrapfig, tabularx, multirow, makecell, subcaption}
\usepackage[flushleft]{threeparttable}
\usepackage[table]{xcolor} %einfärben von tabellenreihen
%\usepackage[round, authoryear]{natbib}
\usepackage[
backend=biber,
style=apa,
maxbibnames=6,
minbibnames=3,
natbib=true          % <- erlaubt weiterhin \citep{}, \citet{} wie gewohnt
]{biblatex}
\addbibresource{References.bib}


\usepackage{chemformula} % Hauptpaket für chemische Formeln
\usepackage{chemfig} % für chemische Strukturformeln

%%%%%%%%%%%%%%%%%%%% für Tabellen/Figurebeschriftung 
\captionsetup{
	labelfont = bf,        % "Figure 1:" fett
	labelsep = colon,       % Doppelpunkt nach der Nummer (Standard, aber explizit)
	font = small,           % gesamte Caption etwas kleiner (optional, oft für Abschlussarbeiten üblich)
	justification = justified   % Blocksatz für die Caption 
}

\setlength{\parindent}{0pt}

\usepackage[hidelinks]{hyperref}

\usepackage[acronym, style = long, toc]{glossaries}

\hypersetup{
     colorlinks   = true, %If links to papers should be colored
     citecolor    = blue, %color of citation text
     urlcolor = black, %if you use an URL, this is how it is colored
     linkcolor = black
}
\usepackage{cleveref}
\graphicspath{{figures/}}

%CHANGE SECTION STYLE

\renewcommand{\abstractname}{\Large ABSTRACT}
% \renewcommand{\abstractnamefont}{\normalfont\Large\bfseries}
\renewcommand{\listfigurename}{LIST OF FIGURES} % Uppercase List of Figures
\renewcommand{\listtablename}{LIST OF TABLES} % Uppercase List of Figures

% larger titles
\titleformat{\section}{\normalfont\Large\bfseries}{\thesection}{0.5em}{} % Larger font for section titles
\titleformat{\subsection}{\normalfont\large\itshape}{\thesubsection}{0.5em}{} % Larger font for subsection titles
\titleformat{\subsubsection}{\normalfont\normalsize\itshape}{\thesubsubsection}{0.5em}{} % Larger font for subsubsection titles


%TABLE OF CONTENTS FIXING
\renewcommand*\contentsname{CONTENTS} %Change name
\setcounter{tocdepth}{3} %How many subsections are visible, 1 is only section headings, 3 include subsubsections
